% Fuente íntegra: originales/cuaderno/REORIENTACION_GENEALOGICA_Y_RESULTADOS.md,
% §§2--7. Gestión y procedencia conservadas aparte.
\Needspace{16\baselineskip}
\section{Regularidad genealógica, restitución operatoria y geometría de la acción}
\label{md:restop:seccion}

\subsection{Producción, compatibilidad y restitución}
\label{md:restop:jerarquia}

El fenómeno primario es una dinámica que conserva procedencia al producir
nuevas observaciones. Se reúnen tres operaciones con funciones matemáticas
diferenciadas:
\begin{enumerate}
 \item \emph{Producción.} La actualización APP--TRIT--TPK prolonga el estado
 y sus regiones correlacionadas. La fase puede retornar mientras la memoria
 avanza.
 \item \emph{Compatibilidad.} Las lecturas regionales, angulares, de acción
 y de incidencia obedecen relaciones comunes. Su variación debe respetar
 esas relaciones para corresponder a un mismo estado.
 \item \emph{Restitución.} Conservar la información complementaria permite
 recuperar valores anteriores, productos de operadores y diferencias
 entre órdenes de composición.
\end{enumerate}
La comparación entre procedimientos para obtener $\pi$ debe atender a esos
tres niveles. La tesis específica de HMT reside en la construcción conjunta
y sus relaciones, además de la producción de expansiones numéricas.
Una coincidencia decimal o la periodicidad de una fase no describen por sí
solas el alcance de la construcción. La distinción identifica el objeto
matemático pertinente para una comparación histórica; el desarrollo focal
no establece, por sí mismo, una prioridad mundial.

\subsection{Regímenes de vacancia y retorno de la fase de capacidad}
\label{md:restop:vacancias}

\subsubsection{Calendario de capacidad y selector trítico}
\label{md:restop:antecedentes-capacidad}

El lector de capacidad recibe la compactificación interna $729/1000$ y
distingue el índice de refinamiento del índice de capacidad. Con la
convención de la sección~\ref{vac-capacidad-trit},
\begin{equation}
 \delta=\log_{10}(10/9),\qquad
 p_j=\left\lfloor\frac{j}{\delta}\right\rfloor,\qquad
 v_j=p_j-j,\qquad j\geq1.
 \label{md:restop:capacidad}
\end{equation}
Los intervalos entre eventos son
\begin{equation}
 h_j=p_{j+1}-p_j\in\{21,22\},\qquad
 s_j=22-h_j.
 \label{md:restop:intervalos}
\end{equation}
La clase trítica obedece
\begin{equation}
 [v_{j+1}]_3=[v_j-s_j]_3.
 \label{md:restop:clase-tritica}
\end{equation}
Un intervalo corto cambia la clase; uno largo la conserva.
Las separaciones $6/7$ describen las mesetas de este régimen inducido.
La construcción mantiene su tipo de factor de observación del estado TPK;
ni sustituye dicho estado ni identifica esas separaciones con signos de
carga o con cifras de $\pi$ o de $e$.
El selector canónico es
$M_{\mathrm{ph}}(1+[v_j]_9)$, cuya pauta depende de la clase módulo tres.

\subsubsection{Restitución del régimen local y desplazamiento módulo nueve}
\label{md:restop:consecuencia-capacidad}

Sea
\begin{equation}
 \beta_{\mathrm{cap}}=22-\delta^{-1},\qquad
 q_k=\left\lfloor\frac{k}{\beta_{\mathrm{cap}}}\right\rfloor,\qquad k\geq1.
 \label{md:restop:pendiente-inducida}
\end{equation}
Las posiciones $q_k$ son los índices de los intervalos cortos. En efecto,
\begin{equation}
 p_j=22j-\lceil j\beta_{\mathrm{cap}}\rceil,\qquad
 s_j=\lceil(j+1)\beta_{\mathrm{cap}}\rceil
     -\lceil j\beta_{\mathrm{cap}}\rceil.
 \label{md:restop:eventos-cortos}
\end{equation}
La irracionalidad de $\delta$ se obtiene de las valuaciones de $2,3,5$
en la igualdad que impondría $(10/9)^q=10^p$; implica la de
$\beta_{\mathrm{cap}}$. Por ello, los índices positivos no presentan
ambigüedad de extremos. La discontinuidad correspondiente al entero $k$
ocurre en $j=\lfloor k/\beta_{\mathrm{cap}}\rfloor$.

\begin{proposition}
\label{md:restop:proposicion-regimen}
Entre los eventos de índices $q_k$ y $q_{k+3}$, tomando los intervalos con
índices en $[q_k,q_{k+3})$, retorna la clase TRIT, mientras cambia la fase
de capacidad módulo nueve.
\end{proposition}
\begin{proof}
Hay exactamente tres intervalos cortos en ese tramo.
Las cotas $1/7<\beta_{\mathrm{cap}}<3/20$ dan
$20<3/\beta_{\mathrm{cap}}<21$, por lo que
\begin{equation}
 m_k=q_{k+3}-q_k\in\{20,21\}.
 \label{md:restop:duracion-tritica}
\end{equation}
La suma de longitudes y la ganancia de capacidad son, respectivamente,
\begin{equation}
 \Delta p=22m_k-3\in\{437,459\},\qquad
 \Delta v=21m_k-3\in\{417,438\}.
 \label{md:restop:ganancias}
\end{equation}
En ambos casos $\Delta v\equiv0\pmod3$. En cambio,
\begin{equation}
 \boxed{\Delta v\equiv3\ \text{o}\ 6\pmod9.}
 \label{md:restop:desplazamiento-nonadico}
\end{equation}
Queda probado el enunciado.
Las cotas de $\beta_{\mathrm{cap}}$ disponen de certificados enteros:
\begin{equation}
 \begin{aligned}
 10^{146}>9^{153}&\ \Longrightarrow\
 \delta>\frac7{153},\\
 10^{417}<9^{437}&\ \Longrightarrow\
 \delta<\frac{20}{437}.
 \end{aligned}
 \label{md:restop:certificados-cotas}
\end{equation}
La monotonía de $22-1/x$ produce
$1/7<\beta_{\mathrm{cap}}<3/20$. Así, las cotas se apoyan en
desigualdades enteras y no en una sustitución decimal redondeada.
\end{proof}

El valor del selector
$\tau_j=M_{\mathrm{ph}}(1+[v_j]_9)$ retorna porque depende de la clase
módulo tres. Este retorno conserva su distinción con el estado TRIT
completo y con una vuelta cronológica de $\Gamma_9$.
Ambas duraciones aparecen infinitas veces por la densidad de la rotación
irracional de pendiente $1/\beta_{\mathrm{cap}}$.

La regularidad tiene niveles distintos: se restaura el régimen local
mientras sigue desplazada la fase de capacidad. El retorno nonádico del
estado tiene, además, su incremento propio de memoria. Cada objeto
conserva, por tanto, su noción de retorno y su duración.
Los números $437$ y $459$ son duraciones de esta observación inducida,
sin atribuirles el estatuto de constantes universales adicionales.
El resultado reúne explícitamente las fórmulas previas; su presentación
como consecuencia focal no constituye una afirmación de prioridad
histórica frente a desarrollos no comparados exhaustivamente.

\subsection{Participación de la memoria en la composición operatoria}
\label{md:restop:memoria}

\subsubsection{Representación ampliada de lectura y memoria}
\label{md:restop:estructura-partida}

La representación de lectura y memoria se obtiene separando ocho
componentes y una novena. En una misma fibra, con transporte de referencia
identidad, el transporte interno $C$ induce
\begin{equation}
 T_C=\frac{8I+C}{9},\qquad
 D_C=\frac{\sqrt8(C-I)}{9},\qquad
 R_C=\frac{I+8C}{9},
 \label{md:restop:bloques}
\end{equation}
\begin{equation}
 \mathcal U_C=
 \begin{pmatrix}
 T_C&D_C\\D_C&R_C
 \end{pmatrix}.
 \label{md:restop:transporte-ampliado}
\end{equation}
El acoplamiento procede del cambio de base de la descomposición nonádica.
Para transportes que conservan la forma declarada, conserva la forma
ampliada. La identidad de composición es
\begin{equation}
 \mathcal U_{C_2}\mathcal U_{C_1}
 =\mathcal U_{C_2C_1},\qquad
 T_2T_1+D_2D_1=\frac{8I+C_2C_1}{9},
 \label{md:restop:composicion-ampliada}
\end{equation}
con $T_i=T_{C_i}$ y $D_i=D_{C_i}$.
El segundo término corresponde a la memoria producida en la primera
operación y reintroducida en la segunda. Apartarla cambia el
procedimiento: produce $T_2T_1$ en lugar de la lectura de la composición
ampliada.

\subsubsection{Reparto de la sensibilidad al orden de composición}
\label{md:restop:conmutadores}

\begin{proposition}
\label{md:restop:proposicion-orden}
Con la convención $[A,B]=AB-BA$,
\begin{equation}
 [T_2,T_1]=\frac1{81}[C_2,C_1],\qquad
 [D_2,D_1]=\frac8{81}[C_2,C_1].
 \label{md:restop:reparto-conmutadores}
\end{equation}
En consecuencia, la diferencia de lectura entre ambos órdenes, con la
memoria reintroducida, es
\begin{equation}
 \boxed{(T_2T_1+D_2D_1)-(T_1T_2+D_1D_2)
 =\frac19[C_2,C_1].}
 \label{md:restop:orden-completo}
\end{equation}
\end{proposition}
\begin{proof}
Al expandir los productos de $T$, los términos constantes y lineales se
cancelan en la resta; queda el conmutador dividido por $81$.
Al expandir los productos de $D$ se produce la misma cancelación y el
factor es $8/81$. Su suma es $1/9$.
\end{proof}

El mecanismo de reentrada forma parte de la representación ampliada.
La comparación de los dos órdenes hace explícita una consecuencia:
en esta representación, la contribución de memoria a la respuesta
antisimétrica es ocho veces la retenida por la composición de lecturas
aisladas. La respuesta completa es nueve veces esta última. Las
proporciones corresponden a operadores de respuesta, no a porcentajes
de energía.

La restitución conserva la no conmutatividad y su intensidad de lectura.
En este caso, la compresión aislada atenúa un conmutador no nulo sin
anularlo. Una aplicación a fuerzas físicas conserva además los operadores
de interacción y el funcional de trabajo correspondiente; la identidad
algebraica no los identifica por sí sola.

\subsubsection{Contraste operatorio y normalización nonádica}
\label{md:restop:contraste}

Para el patrón algebraico general $(n-1)+1$, las mismas cuentas dan los
factores $1/n^2$, $(n-1)/n^2$ y $1/n$. La arquitectura nonádica fija
$n=9$. Se distinguen así la ley de composición y el valor concreto de
su normalización.

Cuando las operaciones efectivas de una aplicación se representan por
$C_1,C_2$, los procedimientos que apartan o reintroducen la memoria
predicen respuestas de orden distintas. El contraste recae en la
composición operacional y en la accesibilidad de la memoria, conservando
los coeficientes de la representación.

\subsubsection{Normas de lectura y memoria en la inversión central}
\label{md:restop:treinta-y-dos}

Para $C=-I$, se obtiene
\begin{equation}
 T^*T=\frac{49}{81}I,\qquad
 D^*D=\frac{32}{81}I.
 \label{md:restop:normas}
\end{equation}
El numerador $32$ procede de $8|{-1}-1|^2$.
En las nueve componentes de la memoria, la norma incluye
$8^2+8=72$, con denominador $27^2$. El cuadrado $64$ participa en
ese cálculo junto con las otras ocho contribuciones.
Una identificación con una región de $64$ casillas exige el mapa de
incidencia de la región y no se obtiene duplicando el numerador $32$.
La lectura geométrica de un soporte $8\times8$ frente a otro
$9\times9$ conserva, por tanto, su problema de conexión con la
aplicación concreta; la igualdad numérica aislada no lo sustituye.

\clearpage
\subsection{Carácter de acción y coordenada completa de reciprocidad}
\label{md:restop:h5}

\subsubsection{Coeficientes genealógicos y continuación regional}
\label{md:restop:coeficientes}

La matriz central y sus invariantes son
\begin{equation}
 B_c=\begin{pmatrix}7&2\\2&7\end{pmatrix},\qquad
 \lambda_+=9,\qquad\lambda_-=5.
 \label{md:restop:matriz-central}
\end{equation}
Intervienen además la órbita de longitud cuatro del paso tres, el
calendario de longitud doce y el cociente censal $104976/1944=54$.
Las relaciones de coeficientes son
\begin{equation}
 \begin{aligned}
 \frac9{16}&=\frac{\lambda_+}{4^2},&
 \frac59&=\frac{\lambda_-}{\lambda_+},\\
 \frac7{48}&=\frac{\operatorname{tr}B_c/2}{4\cdot12},&
 \frac1{54}&=\frac{1944}{104976}.
 \end{aligned}
 \label{md:restop:coeficientes-h5}
\end{equation}
La asignación de grados y signos pertenece a la definición del carácter;
sus cardinales aislados no la determinan. El lector completo es
\begin{equation}
 H_5=\frac12\sqrt{\frac{1000\alpha}{\varphi}}-\alpha
 +\frac9{16}\alpha^2-\frac59\alpha^3
 +\frac7{48}\alpha^4-\frac1{54}\alpha^5,
 \label{md:restop:h5-completo}
\end{equation}
\begin{equation}
 \begin{aligned}
 D_A&=\exp(-100\pi\alpha/9),\\
 \mathcal C_\pi&=169\alpha^6+
 \frac{D_A\alpha^7}{1-D_A\alpha},\\
 \eta_{\mathrm{ret}}&=H_5-\frac{90}{\pi}\mathcal C_\pi.
 \end{aligned}
 \label{md:restop:retorno-regional}
\end{equation}
La fracción conserva la continuación infinita definida por
\begin{equation}
 \mathscr R(z)=D_Az^7+D_Az\mathscr R(z),
 \label{md:restop:renovacion}
\end{equation}
con peso inicial unitario. La continuación permanece íntegra en el
lector; un polinomio truncado no sustituye ese objeto. La construcción
de los coeficientes y la renovación se reúnen en la
sección~\ref{sec:revision-planck}.

\subsubsection{Clasificación de la reciprocidad radial orientada}
\label{md:restop:geometria}

Se conservan $A_{\deg}=1000\alpha$,
$C^*_{\deg}=2(\eta_{\mathrm{ret}}+\alpha)$ y el recuperador radial del
teorema~\ref{rec:thm:accion}. Para mantener ambas orientaciones, sean
\begin{equation}
 \epsilon\in\{+1,-1\},\qquad
 q_{\mathrm{rad}}=R_\star^4>0,\qquad
 \Xi_{\mathrm{rad}}=
 \epsilon\frac{1-q_{\mathrm{rad}}}{1+q_{\mathrm{rad}}}.
 \label{md:restop:par-radial}
\end{equation}
El recuperador \eqref{rec:eq:eta-orientada} se escribe
\begin{equation}
 \frac{\eta_{\mathrm{ret}}}{\alpha}=500\Xi_{\mathrm{rad}}-1.
 \label{md:restop:recuperador}
\end{equation}
La involución
$(\epsilon,q_{\mathrm{rad}})
 \mapsto(-\epsilon,q_{\mathrm{rad}}^{-1})$
conserva esa acción. La coordenada que reúne esta estructura es
\begin{equation}
 \boxed{\chi_{\mathrm{rad}}=\epsilon\log q_{\mathrm{rad}}.}
 \label{md:restop:chi}
\end{equation}

\begin{proposition}
\label{md:restop:proposicion-orbitas}
$\chi_{\mathrm{rad}}$ clasifica exactamente las órbitas de esa involución
sobre $\{\pm1\}\times\mathbb R_{>0}$.
\end{proposition}
\begin{proof}
Cambiar simultáneamente el signo e invertir $q_{\mathrm{rad}}$ conserva
$\chi_{\mathrm{rad}}$. Si dos pares tienen el mismo valor de
$\chi_{\mathrm{rad}}$, con igual signo tienen el mismo
$q_{\mathrm{rad}}$; con signos opuestos, sus radios cuárticos son
recíprocos. Son exactamente los dos casos de pertenencia a una misma
órbita. El argumento incluye $q_{\mathrm{rad}}=1$, cuya órbita conserva
las dos orientaciones.
\end{proof}

La identidad
$(1-q)/(1+q)=-\tanh(\log q/2)$ da
\begin{equation}
 \Xi_{\mathrm{rad}}
 =-\tanh(\chi_{\mathrm{rad}}/2),\qquad
 \boxed{\eta_{\mathrm{ret}}
 =\alpha[-500\tanh(\chi_{\mathrm{rad}}/2)-1].}
 \label{md:restop:accion-chi}
\end{equation}
La compatibilidad con el carácter completo adopta la forma
\begin{equation}
 \boxed{\chi_{\mathrm{rad}}
 =-2\operatorname{artanh}
 \left(\frac{H_5+\alpha-(90/\pi)\mathcal C_\pi}{500\alpha}\right).}
 \label{md:restop:chi-caracter}
\end{equation}
El término $-\alpha$ de $H_5$ se cancela con el $+\alpha$ del
contraste angular antes de esta lectura. El denominador de la
continuación regional permanece completo. La igualdad recibe las
salidas y los lectores HMT anteriores, y reúne sus dos representaciones;
no define retroactivamente $\alpha$.

En la rama positiva de acción y radio finito,
\begin{equation}
 0<\eta_{\mathrm{ret}}<499\alpha,\qquad
 \chi_{\mathrm{rad}}
 <-2\operatorname{artanh}(1/500).
 \label{md:restop:dominio-positivo}
\end{equation}
Con $\alpha$ fija,
\begin{equation}
 \frac{\partial\eta_{\mathrm{ret}}}{\partial\chi_{\mathrm{rad}}}
 =-250\alpha\,\operatorname{sech}^2(\chi_{\mathrm{rad}}/2)<0.
 \label{md:restop:derivada-accion}
\end{equation}
Por tanto, la sección adimensional de acción determina unívocamente la
clase radial orientada en esta representación, y recíprocamente.
Para variaciones conjuntas,
\begin{equation}
 d\eta_{\mathrm{ret}}
 =\frac{\eta_{\mathrm{ret}}}{\alpha}\,d\alpha
 -250\alpha\,\operatorname{sech}^2(\chi_{\mathrm{rad}}/2)
 \,d\chi_{\mathrm{rad}}.
 \label{md:restop:diferencial-accion}
\end{equation}

El carácter $H_5$, la corrección regional y la anisotropía radial orientada
mantienen así una relación de compatibilidad y no constituyen tres
parámetros independientes. El carácter analítico y la geometría de
reciprocidad representan la misma sección de retorno.
Este nivel estructural de la ecuación de Planck permanece explícito
antes de componer la acción con una frecuencia.

La década de acción conserva su genealogía determinantal, independiente
de esta reparametrización:
\begin{equation}
 \hbar_{\mathrm{ret}}
 =10^{-34}\mathcal U_S\eta_{\mathrm{ret}},\qquad
 \hbar_{\mathrm{pre}}=10^{-34}\mathcal U_SH_5.
 \label{md:restop:secciones-accion}
\end{equation}
Se conservan además la acción por vuelta $h=2\pi\hbar$ y la lectura
térmica
\begin{equation}
 k_B\Theta_{\mathrm{clk}}\log3=\frac{h}{T_{108}}
 \label{md:restop:lectura-termica}
\end{equation}
en su ámbito. El reloj físico de una realización conserva su constructor:
la coordenada $\chi_{\mathrm{rad}}$ clasifica esta reciprocidad y no la
totalidad de historias TPK. La clasificación por
$\chi_{\mathrm{rad}}$ y la relación diferencial reúnen las representaciones
existentes de la acción; mantienen su carácter de formalización focal y no
constituyen una segunda derivación independiente de Planck.

\subsection{Circuitos de fase, memoria y conservatividad}
\label{md:restop:conservatividad}

La holonomía nonádica restaura la fase e incrementa la memoria, según
la sección~\ref{mono-ampl:vueltas}. Si un recorrido cerrado en las fases
aumenta el contador entero, ese incremento no puede escribirse como
diferencia de un potencial dependiente únicamente de la fase.

La demostración se obtiene por telescopía: la suma
\begin{equation}
 \sum_j\bigl[F(g_{j+1})-F(g_j)\bigr]=0
 \label{md:restop:potencial-fase}
\end{equation}
en un circuito de fases contrasta con el incremento de memoria, igual
a uno por vuelta. Sobre el estado ampliado, el contador $m$ registra
ese incremento. El recorrido cerrado en la proyección conservaba,
por tanto, una apertura en el estado completo.

El estudio de conservatividad comienza por declarar el espacio en el
que se cierra el recorrido y las coordenadas que participan en la
dinámica. No conmutatividad y no conservatividad son propiedades
diferentes. La sección~\ref{md:restop:memoria} controla el orden de
operaciones; esta sección distingue un circuito proyectado de un circuito
completo. Para estudiar una fuerza concreta, el funcional de acción o
trabajo debe actuar sobre el mismo espacio de estados y el mismo transporte.

El crecimiento de memoria no asigna por sí mismo una producción
termodinámica. La sección~\ref{md:cron:crecimiento} demuestra que la
información de bloques puede crecer con tasa entrópica nula.
La restitución operatoria añade una precisión: conservar información
puede significar conservar qué transformaciones se pueden componer,
además de mantener los símbolos almacenados.

\subsection{Compatibilidad dinámica entre observaciones}
\label{md:restop:sintesis}

La unidad de estudio es la compatibilidad dinámica entre observaciones,
con la memoria como parte de la operación. El valor de una constante
constituye una lectura; su estructura contiene restricciones sobre otras
lecturas y sobre las transformaciones admisibles. La prolongación del
análisis consiste en obtener relaciones que impidan modificar de manera
independiente valor, orientación, memoria, forma y acción.

Las tres elaboraciones focales recorren ese mismo eje:
\begin{itemize}
 \item Un ciclo completo de regímenes TRIT conserva todavía un
 desplazamiento nonádico de capacidad: la regularidad local no agota
 el retorno.
 \item La reentrada de memoria modifica exactamente la respuesta al orden
 de dos operaciones: el historial interviene causalmente en la composición.
 \item El carácter $H_5$ y la geometría radial orientada determinan
 mutuamente su sección de retorno: la acción admite una representación
 estructural recuperable.
\end{itemize}
Estas consecuencias continúan la genealogía de $\pi,\varphi,e,\alpha$ y
explicitan qué significa conservarla al pasar a las operaciones y la
acción. La composición posterior reúne esos niveles mediante los mapas
comunes ya construidos, manteniendo la distinción entre sus índices y
sus memorias aunque compartan una denominación.
